\chapter*{Agradecimientos}

Expreso mi profundo agradecimiento a mi familia por su apoyo constante, su paciencia y la confianza que me brindaron durante todo el proceso de desarrollo de esta tesis. Su acompañamiento fue fundamental para sostener el esfuerzo académico y personal que implicó llevar esta propuesta a un nivel de formulación técnica y aplicada.

Agradezco de manera especial a mi asesor, MSc. MBA. Ing. René Mitma Ramírez, por su orientación, observaciones y acompañamiento a lo largo de la elaboración de este trabajo. Sus recomendaciones permitieron fortalecer el enfoque metodológico, la estructura del documento y la pertinencia de la propuesta planteada.

Asimismo, agradezco a la Universidad Nacional de Ingeniería y a la Facultad de Ingeniería Eléctrica y Electrónica por la formación recibida, así como por brindar el entorno académico necesario para el desarrollo de esta investigación. Del mismo modo, reconozco el valor de las fuentes institucionales, técnicas y científicas consultadas, las cuales hicieron posible sustentar esta propuesta de manera rigurosa.

Finalmente, agradezco a todas las personas y profesionales del ámbito cultural y tecnológico que, directa o indirectamente, inspiran la búsqueda de soluciones innovadoras para mejorar la accesibilidad, la mediación cultural y la experiencia de visita en museos y espacios patrimoniales del Perú.
